\section{Concurrency and Parallelelisation}
\textcolor{Blue}{Parallelism}: multiple subtasks run at the same time, with a
processor that runs multiple instructions, or with multiple processors.\\
\textcolor{Blue}{Concurrency}: the ``illusion'' of parallelism.
\begin{center}
  $Parallel\Rightarrow Concurrent,Concurrent\nLeftarrow Parallel$
\end{center}
Parallelisation is \textcolor{Bittersweet}{non-deterministic}: no order of
which processes start and complete.
\paragraph{Parallelisability}
\begin{itemize}
  \item No Interference: No modification of the stream during execution of
    terminal operation. Throws
    \textcolor{Bittersweet}{ConcurrentModificationException}
  \item Stateless: The result does not depend on any state that might change
    during executiion of the stream \textcolor{ForestGreen}{Previous elements,
    Input}
  \item No Side-effects: No concurrent modification of non-thread safe objects
    that may result in incorrect behaviour
\end{itemize}
Operations that fulfil these conditions are known as \textcolor{Blue}{embarrassingly parallel}.
\begin{center}
  Parallelising $\neq$ Performance
\end{center}
Creating too many threads will \textcolor{Bittersweet}{outweigh benefits} of
parallelisation.
For \textbf{streams}, particularly for
\begin{verbatim}
stream.reduce(U i,
   BiFunc<U, ? super T, U> acc,
   BinOp<U> comb>)
\end{verbatim}
which accumulates each substream and combines the result, parallelisation
further requires:
\begin{itemize}
  \item $\verb|combiner.apply(i, x)|\equiv\verb|x|,\forall x$
  \item \verb|combiner| and \verb|accumulator| are associative.
  \item \verb|combiner| and \verb|accumulator| are compatible.\\
  $\verb|comb.apply(u, acc.apply(i, v))|\equiv\\\verb|acc.apply(u,v)|,\forall u,v$
\end{itemize}
Streams themselves may define encounter order:
\begin{multicols*}{2}
  \textcolor{ForestGreen}{\textbf{Ordered}}
  \begin{itemize}
    \item \verb$Stream::iterate$
    \item \verb$Stream::of$
    \item Ordered Collections
  \end{itemize}
  \textcolor{Bittersweet}{\textbf{Unordered}}
  \begin{itemize}
    \item \verb$Stream::generate$
    \item Unordered Collections
  \end{itemize}
\end{multicols*}
Operations that respect encounter order: \verb|distinct|, \verb|sorted|,
\verb$findFirst$, \verb$limit$, \verb$skip$ are expensive.
\verb|stream.unordered()| if original order is not important.
\paragraph{Threads}
In Java, threads divide computation into subtasks, to improve utilisation of
processor.
\begin{center}
  \verb$new Thread(RUNNABLE).start();$
\end{center}
The execution flow of each thread can be separately programmed with
\verb$Runnable.run()$. Program exits only after \emph{ALL} created threads run
to their completion.
\paragraph{CompletableFuture<T>}
\emph{A MONAD}. A \verb$CF$ is either completed or not complete.\\
To \textbf{Create a CF}:
\begin{itemize}
  \item \verb|CF<U> CF.completedFuture(U value)|
  \item \verb|CF<Void> runAsync(Runnable r)|
  \item \verb|CF<T> supplyAsync(Supplier<T> s)|
\end{itemize}
To \textbf{Chain CFs}:
\begin{itemize}
  \item \verb|CF<Void> thenAccept(Consumer<T> c)|
  \item \verb|CF<U> thenApply(Function<T, U> map)|
  \item \verb|CF<U> thenCompose(Function<T, CF<U>> map)|
  \item \begin{verbatim}
CF<U> thenCombine(CF<V> other,
                  BiFunc<T, V, U> fn)
    \end{verbatim}
  \item \verb|CF<Void> thenRun(Runnable r)|
\end{itemize}
To \textbf{get CF result}:
\begin{itemize}
  \item \verb|T get()| throws \textcolor{Bittersweet}{InterruptedException, ExecutionException}
  \item \verb|T join()| throws any \textcolor{Bittersweet}{Exception} that occurs
\end{itemize}
\textbf{Handling Exceptions}:
\begin{itemize}
  \item \verb|CF<T> exceptionally(Func<Throwable, T> f)|
  \item \verb|whenComplete(BiCons<T, Throwable> cons)|
  \item \verb|CF<R> handle(BiFunc<V, E, R> f)|
\end{itemize}
\begin{center}
  \vspace{-4mm}
  \includegraphics[width=\linewidth]{recursivetask}
  \vspace{-5mm}
\end{center}
\begin{enumerate}
  \item Idle Thread: Take head of own deque. If empty, take tail of another deque, and compute.
  \item \verb|fork()| called: Invoker adds itself to the head of current deque.
  \item \verb|join()| called: If complete, read result and join returns. If not
    stolen, compute on current thread. If stolen, do another task while waiting.
\end{enumerate}
Order of fork join: ``Palindrome'', with only ONE compute in the middle
